\scanpage{27}

% L1–2
\uline{Th.} Soit 1) $\sigma$ un ensemble de couples $(S,\lambda)$, $S$ catégorie, $\lambda=(\lambda',\lambda'')$ ens. de cônes (proj. et ind.) dans $S$

% L3–6
2) $\rho$ un ens de \deleted{diagrammes} \inserted{systèmes} $(\deleted{\unsure{R}},R_1,\delta_1,R_2,\delta_2,\varphi\deleted{\unsure{,R_3}})$ où $(R_i,\delta_i)$ \inserted{$(i=1,2)$} est une catégorie avec ens $\delta_i=(\delta'_i,\delta''_i)$ de cônes. $R$ une catégorie, \deleted{\unsure{$R_1,R_2$} des foncteurs \unsure{$\varphi_1$}} $\varphi$ un foncteur $R_2\mathrel{\overset{\varphi}{-}}R_1$ tel
% L5, right, cancelled triangle
\diagram{\begin{tikzcd}[ampersand replacement=\&,row sep=small,column sep=small,execute at end picture={\draw (\tikzcdmatrixname-2-1.south west)--(\tikzcdmatrixname-1-2.north east) (\tikzcdmatrixname-2-2.south west)--(\tikzcdmatrixname-1-3.north east);}]
{}\&\deleted{\unsure{R}}\arrow[dl,"\deleted{\varphi_1}"']\arrow[dr,"\deleted{\varphi_2}"]\&{}\\
\deleted{R_1}\&{}\&\deleted{R_2}
\end{tikzcd}}
que $\varphi(\delta_2)\subset\delta_1$ \deleted{$(R_1,\delta_2)$}

% L7
Soit $\Delta=\Delta_{\sigma,\rho}$ la \inserted{sous 2-}catégorie de Cat définie ainsi

% L8
a) Les objets sont les catégories $E$ ayant les propriétés suivantes
% L9–10
$a_1)$ $\forall(S,\lambda)\in\sigma$ et tout foncteur $S\xrightarrow{u}E$, les systèmes projectifs et inductifs correspondant aux cônes $\in u(\lambda)$ ont une limite
% L11–13
$a_2)$ $\forall(\deleted{\unsure{R}}R_1,\delta_1,\ldots)\in\rho$, \deleted{et tout foncteur $u_1:R_1-E$ qui est $\delta_1$-exact, il existe un foncteur $u_2:R_2-E$ qui est $\delta_2$-exact et tel que $u_1\varphi_1\simeq u_2\varphi_2$} \inserted{$u\mapsto u\circ\varphi$ induit une équivalence}
$\underline{\mathrm{Hom}}_{\delta_1}(R_1,E)\xrightarrow{\sim}\underline{\mathrm{Hom}}_{\delta_2}(R_2,E)$

% L14–19
b) Pour $E,E'\in\operatorname{Ob}\Delta$, $\underline{\mathrm{Hom}}_\Delta(E,E')$ est la sous-catégorie \inserted{(\unsure{str.|strte|stab.})} pleine de $\underline{\mathrm{Hom}}_{\mathrm{Cat}}(E,\unsure{E|E'|F})$ formée des foncteurs $F:E\to E'$ tels que l'on ait : $\forall(S,\lambda)\in\sigma$ et tout foncteur $u:S-E$ qui est $\lambda$-exact, $F\circ u$ est $\lambda$-exact. (Il est clair que cette condition est stable par isomorphisme et par composition, et vérifiée \deleted{par les} \inserted{si $F$ est une} équivalence de \unsure{catégorie|catégories}.)

% L20
On suppose satisfaites les conditions
% L21–25
A) $\sigma$ et $\rho$ sont petits \deleted{les} $\forall\,\unsure{(S,\lambda)|S,\lambda)}\in\sigma$ $S$ est petit et $\lambda$ est petit, et les \deleted{diagrammes} \inserted{catégories} d'indices des cônes $\in\lambda$ sont petites, de même pour tout $(\deleted{\unsure{R}},R_1,\delta_1,R_2,\delta_2,\varphi\deleted{\unsure{\varphi_2}})\in\rho$ les catégories \deleted{\unsure{$R_2$,}} $R_i$ $(i=1,2)$ sont petites, les $\delta_i$ $(i=1,2)$ sont \unsure{petit|petits} et les catégories d'indices des cônes $\in\delta_i$ sont petites.

% L26–27
B) $\forall(R_1,\ldots)\in\rho$, $(R_{\unsure{i|1}},\delta_{\unsure{i|1}})\in\sigma$ pour $i=1,2$. \inserted{(ou du moins : la condition d'exactitude des cat et \deleted{\illegible{1}} \unsure{des|de|le} foncteurs définie par $(R_i,\delta_i)$ est conséquence de celle définie par $\sigma$)} [On ne veut introduire des conditions d'exactitude supplémentaires b)
% archive sheet number: 26
